\chapter{OCR --- Reconnaissance Optique de Caractères}

\section{Identification}

\begin{center}
\begin{tabular}{ll}
\toprule
\textbf{Champ} & \textbf{Valeur} \\
\midrule
Modèle actif    & QwenVL 2.5-7B (vLLM) \\
Statut scoring  & À implémenter \\
Fichier         & \texttt{DataPipelines/VLMModel.py} \\
\bottomrule
\end{tabular}
\end{center}

\section{Fonctionnalité}

Le composant OCR convertit les images en texte exploitable. Il est déclenché
sur chaque n\oe ud image portant le flag \texttt{OCR\_FLAG}. Le texte produit
est découpé en chunks, vectorisé par BGE-M3, traduit si nécessaire, puis
indexé pour la recherche.

\textbf{Entrée.} Image binaire (JPEG, PNG, BMP) $> 100 \times 100$ px, ratio
max/min $< 10$.\\
\textbf{Sortie.} Texte brut multi-lignes, ordre de lecture reconstitué.

\section{Modèles évalués}

\begin{center}
\begin{tabular}{llll}
\toprule
\textbf{Modèle} & \textbf{Version} & \textbf{Benchmarké} & \textbf{Retenu} \\
\midrule
EasyOCR + Tesseract & 1.7 + 5.x & Non (legacy)   & Non \\
QwenVL 2.5-7B       & 7B Q8     & Non            & Oui \\
Marker              & TDB       & Non        & Non \\
Docling (IBM)       & TDB       & Non        & Non \\
InternVL 3.5-7B     & TDB       & Non        & Non \\
\bottomrule
\end{tabular}
\end{center}

\section{Protocole d'évaluation}

\noindent\textbf{Métriques.} TBD.

\noindent\textbf{Dataset à constituer.}  TBD 

\section{Fonction de scoring de confiance}

\subsection{Mode legacy --- EasyOCR + Tesseract}

%Soit $W^{+} = \{w_i : c_i \geq 0\}$ les mots détectés avec score Tesseract
%non nul.

% \begin{equation}
%   s_{\text{legacy}}(x, \hat{y})
%   = \frac{1}{|W^{+}|} \sum_{i \in W^{+}} \frac{c_i}{100}
%   \;\in [0, 1]
%   \label{eq:ocr_legacy}
% \end{equation}

% La rotation de l'image retient la transcription maximisant ce score.
TBD
\subsection{Mode VLM / QwenVL (actif) --- Proposition}

% Le VLM ne produit pas de scores par token via l'API actuelle. Deux approches
% sont proposées.

% \paragraph{Option A --- Confiance verbalisée (recommandée).}
% Ajout d'une instruction de post-génération~:
% \textit{\og Transcris le texte. Termine par : CONFIANCE: <entier 0--100>.\fg{}}

% \begin{equation}
%   s_{\text{VLM}}(x, \hat{y})
%   = \frac{c_{\text{verb}}}{100} \cdot \delta_{\text{format}}
%   \;\in [0, 1]
%   \label{eq:ocr_vlm}
% \end{equation}

% où $\delta_{\text{format}} = 1$ si le token CONFIANCE est présent et
% parseable, $0.5$ sinon.

% \paragraph{Option B --- Self-Consistency ($K$ inférences).}

% \begin{equation}
%   s_{\text{SC}}(x)
%   = 1 - \frac{1}{K(K-1)} \sum_{i \neq j} \mathrm{NED}(\hat{y}_i, \hat{y}_j)
%   \;\in [0, 1]
%   \label{eq:ocr_sc}
% \end{equation}

% Coût : $K \times$ latence.
TBD
\paragraph{Seuils provisoires.}

\begin{center}
\begin{tabular}{lll}
\toprule
\textbf{Plage} & \textbf{Interprétation} & \textbf{Action} \\
\midrule
TBD           & Confiance moyenne & Relecture recommandée \\
TBD           & Haute confiance  & Exploitable \\
TBD           & Faible confiance  & Relecture obligatoire \\
\bottomrule
\end{tabular}
\end{center}

\section{Justification scientifique}

\section{Validation empirique}

